% ============================================================
% Subgoal Scale & Reality-Verified Speculative Consumption
% Preliminary internal draft -- PushT results only, cross-environment
% generalization pending. Structured to convert to a conference template
% (ICLR/NeurIPS-style) later; using a plain article class for now so it
% compiles standalone without an external .sty (the ICLR template folder
% is not vendored in this repo -- fetch github.com/ICLR/Master-Template
% and swap the preamble below for the real submission pass).
% ============================================================
\documentclass[11pt]{article}
\usepackage[margin=1in]{geometry}
\usepackage{times}
\usepackage{graphicx}
\usepackage{amsmath,amssymb}
\usepackage{booktabs}
\usepackage{multirow}
\usepackage{tabularx}
\usepackage{enumitem}
\usepackage{subcaption}
\usepackage{xcolor}
\usepackage[colorlinks=true, linkcolor=blue, citecolor=blue]{hyperref}
\usepackage{cleveref}
\usepackage{natbib}

% ---- house macros ----
\newcommand{\ours}[1]{\textbf{#1}}          % highlight our winning cells
\newcommand{\degc}[1]{$#1^\circ$}
\newcommand{\gdm}{\text{GDM}}
% Future-experiment placeholder: renders a red TBD marker.
\newcommand{\TBD}{\textcolor{red}{\textbf{TBD}}}
% A framed red placeholder block standing in for a not-yet-run table.
\newcommand{\tbdblock}[1]{%
  \begin{center}\fcolorbox{red}{red!5}{%
    \parbox{0.9\linewidth}{\centering\color{red}\textbf{[TBD --- #1]}}}%
  \end{center}}

\title{Subgoal Scale and Reality-Verified Speculative Consumption\\
in Hierarchical Latent World-Model Planning\\
\large Preliminary draft --- PushT results only}
\author{}
\date{Draft status: 2026-07-10}

\begin{document}
\maketitle

\begin{center}
\fcolorbox{gray}{gray!8}{\parbox{0.92\linewidth}{\small
\textbf{Status: preliminary draft, PushT-only.} All PushT results below (\Cref{sec:repro}--\Cref{sec:horizon}) are measured and cross-checked against raw logs. \Cref{sec:generalize} (Reacher, TwoRoom, OGBench-Cube) is not yet run --- placeholders only. \Cref{sec:intro} (contributions framing), \Cref{sec:related} (related work), \Cref{sec:method} (formal method description), and \Cref{sec:discussion}--\Cref{sec:conclusion} (discussion/conclusion) are lightly written and will be rewritten once the generalization results land; treat them as scaffolding, not final claims. One open data gap: the random-initialization table (\Cref{sec:randinit}) reports only the seeds we could directly verify against logs --- see the note there.}}
\end{center}

% ============================================================
\begin{abstract}
We study the FF-JEPA hierarchical latent planner --- a frozen LeWM JEPA encoder, a
diffusion subgoal drafter (``GDM''), and a CEM controller --- on PushT, and identify
\emph{subgoal scale}, the spacing $S$ between drafted subgoals with the control loop
matched to it, as the dominant unexplored lever. Success rate over stride is
non-monotone with a clear interior optimum at $S{=}10$: \ours{$97.2\%/96.4\%$}
(\degc{20}/\degc{5} criteria) versus $91.1\%/86.3\%$ for the paper's native $S{=}25$
under an identical protocol. On the paper's own published protocols the winning
configuration \ours{exceeds the published numbers} at both horizons ($96.1{\to}98.4$
at $t{=}25$, $91.8{\to}97.9$ at $t{=}75$, \degc{20}, $n{=}256$); on random
initialization (its hardest setting, published $82.4\%$) it records \ours{zero
failures in every verified episode measured so far} (512/512 across two seeds at
$S{=}10$; a third arm and criterion cell are queued but unconfirmed, see
\Cref{sec:randinit}); and on a fixed-population horizon sweep
the baseline decays monotonically while $S{=}10$ stays flat --- \ours{the gap widens
from $+3.3$ to $+11.9$pp}. A cadence-isolation control shows the mechanism is
\emph{scale matching}, not re-planning frequency. Separately, the one component of
DSpark-style speculative decoding that survives measurement --- a draft--verify--accept
skeleton with the learned verifier replaced by the \emph{achieved environment state}
--- cuts diffusion compute ${\sim}11\times$ versus every-step drafting while remaining
success-rate-neutral at the official criterion. Refinement heads, learned confidence
scheduling, blind multi-step commitment, and a deterministic regressor that dominates
the drafter \emph{offline} each measurably \emph{hurt} in closed loop; we report them
as negative results that causally implicate drafter stochasticity.
\end{abstract}

% ============================================================
\section{Introduction}
\label{sec:intro}

Hierarchical latent world models promise long-horizon planning by decomposing a goal
into a chain of reachable subgoals. FF-JEPA \citep{ffjepa} instantiates this on top of
the LeWM JEPA world model \citep{lewm}: a diffusion model drafts future subgoal latents
and a flat CEM controller \citep{cem} steers toward each in turn, re-planning on a fixed
stride. It reports that hierarchy beats flat CEM on long horizons --- but under an
underspecified pipeline (no released code; unstated diffusion configuration; unstated
evaluation population), and with a subgoal stride fixed at a single value.

We take FF-JEPA as a substrate and ask what actually governs its closed-loop success.
Our contributions:

\begin{enumerate}[itemsep=2pt,topsep=2pt]
\item \textbf{Faithful reproduction and audit} (\Cref{sec:repro}). We rebuild the
pipeline to the episode and iteration, reproduce the paper's headline numbers, and
establish that the benchmark's \emph{coded} success criterion is \degc{20} angular
tolerance (its \degc{5} appears in prose only); we report both throughout.
\item \textbf{The subgoal-scale law} (\Cref{sec:scale}). Retraining the drafter at
$S\in\{5,10,15,25\}$ with the control loop matched to the stride reveals an interior
optimum at $S{=}10$ that beats the paper configuration by $+6.1/+10.1$pp and, on the
paper's own protocols, exceeds its published numbers. A cadence-isolation control
identifies \emph{scale matching}, not re-planning frequency, as the mechanism.
\item \textbf{Reality-verified speculative consumption} (\Cref{sec:spec}). Replacing
DSpark's learned verifier with the achieved environment state yields an
${\sim}11\times$ reduction in diffusion compute at neutral success rate.
\item \textbf{Negative results} (\Cref{sec:neg}, \Cref{sec:regress}). Refinement heads,
blind commitment, learned confidence, and a deterministic regressor all hurt in closed
loop despite matched-or-better \emph{offline} fidelity --- causally implicating the
drafter's conditional \emph{spread} as load-bearing.
\end{enumerate}

% ============================================================
\section{Related Work}
\label{sec:related}
\tbdblock{Related work. Placeholder pointers only --- full positioning pending a
literature-occupancy pass (flagged as an open thread, see \Cref{sec:discussion}).
Axes to cover: (i) hierarchical / subgoal-conditioned latent planning (FF-JEPA
\citep{ffjepa}, LeWM \citep{lewm}, and other JEPA-style world models); (ii)
re-planning cadence and receding-horizon MPC in learned-model control, including
CEM-based planners \citep{cem}; (iii) speculative decoding, both in language
models (the DSpark composition we port, \citep{dspark}) and prior applications
of draft--verify--accept outside text generation; (iv) subgoal/waypoint spacing
and hierarchical RL more broadly; (v) budget-controlled evaluation protocols for
long-horizon manipulation, contrasting our fixed-population sweep with VLWM's
fixed-budget-50 regime \citep{vlwm} and OGBench's benchmark suite \citep{ogbench}.}

% ============================================================
\section{Method}
\label{sec:method}

We modify only three things about the FF-JEPA system (\Cref{sec:repro}): the subgoal
spacing the drafter is trained and queried at, the loop geometry that consumes its
output, and the policy that decides when to re-invoke it. Encoder, drafter
architecture, training objective, and CEM itself are untouched.

\paragraph{Scale-matched retraining.} We retrain the GDM drafter at stride
$S \in \{5,10,15,25\}$ with an identical recipe to the paper's $S{=}25$ configuration
(loss, batch size, learning rate, schedule, epochs --- \Cref{sec:repro}); only the
subgoal spacing used to build training pairs changes, so pair counts are
construction-identical across arms and any SR delta is attributable to spacing alone.
Because the CEM terminal cost is only well-posed when its rollout length is
comparable to the target's distance, the control loop is scaled with the stride:
CEM horizon $=$ replan cadence $=S$ (holding the action-repeat block, sample count,
and top-$k$ fixed). This coupling is definitional, not a second free variable --- see
the cadence-isolation control in \Cref{sec:scale}, which holds targets fixed at
25-step spacing while re-anchoring every 5 steps, and isolates loop cadence from
target spacing as separate mechanisms.

\paragraph{Reality-verified speculative consumption.} \Cref{sec:spec} details the
mechanism; briefly, the drafter proposes an $N$-block once, the policy executes
toward the first position, and after each hop compares the \emph{achieved} latent to
the subgoal just pursued: within tolerance $\tau$, the next drafted position is
consumed with no diffusion call; otherwise the planner re-drafts from the achieved
state. This differs from DSpark's original scheme only in the verifier: DSpark scores
acceptance with a learned confidence head trained on drafter self-consistency; we use
the environment's realized state directly, which is unbiased by construction and
requires no additional training.

\tbdblock{Formal algorithm box / pseudocode for both procedures, and the CEM cost
functional definition, are pending --- currently described in prose above and in
\Cref{sec:scale,sec:spec}.}

% ============================================================
\section{Setting and Faithful Reproduction}
\label{sec:repro}

\paragraph{The FF-JEPA planner.} Observations are encoded by a frozen LeWM JEPA encoder,
$z_t=\mathrm{enc}(o_t)$. A diffusion model $G$ (``GDM'': a 3-block DiT with adaLN-Zero
conditioning, $N{=}3$ future subgoals, $\epsilon$-prediction, trained on
success-filtered expert demonstrations) proposes the next subgoal latent
$\hat z_{sg,m+1}$ from the current encoded observation. CEM then optimizes an action
sequence against a terminal-$L_2$ cost to that subgoal, re-planning every $H$
environment steps. The paper fixes the stride to $H{=}25$.

\paragraph{Reproduction (audit closed).} Our faithful retraining reproduces the paper:
short-horizon $95.31\%$ (paper $96.09\%$; same-set oracle $93.36\%$; $n{=}256$),
long-horizon $t{=}75$ multiseed mean $90.0\%$ (paper $91.80\%$), random-init
$96.5\%/78.1\%$ at \degc{20}/\degc{5} (paper $82.42\%$). Pipeline faithfulness is
over-determined: a from-scratch rebuild of the dense training set reproduces the
original training population to the episode ($9{,}094$) and iteration ($171{,}940$),
and the retrained $S{=}25$ drafter reproduces the reference offline fidelity anchor to
four decimals ($0.1593$ vs.\ $0.1593$ at $m{+}1$).

\paragraph{The \degc{20}/\degc{5} criterion.} FF-JEPA's prose states a \degc{5} angular
tolerance, but the environment's shipped evaluation code enforces \degc{20}; the oracle
(true goal, exact pipeline) reaches only ${\sim}84$--$86\%$ at literal \degc{5} even
under matched-population filtering, making a genuine $96\%$ at \degc{5} for a
\emph{sampled}-subgoal method implausible. We therefore treat \degc{20} as the operative
criterion and \degc{5} as a strict \emph{precision} criterion, reporting both throughout.

\paragraph{Protocols.} Two protocols recur below. \emph{Paper protocol}: $n{=}256$,
paper-matched sampling (reproduction and random-init). \emph{Controlled protocol}:
$n{=}62$ fixed held-out episodes, goal offset $150$, budget $300$ steps, block-only
criterion, CEM seeds 42--45 (curve: 42--49); every arm sees byte-identical episodes and
seeds, so within-table deltas are exact treatment effects. Per-arm
$\mathrm{SE}\approx1.7$pp (4 seeds) / $0.9$pp (8 seeds).

% ============================================================
\section{The Subgoal-Scale Law}
\label{sec:scale}

We retrain the drafter at $S\in\{5,10,15,25\}$ with an identical recipe --- only the
target spacing changes, and pair counts are construction-identical --- and scale the
control loop with the stride (CEM horizon $=$ re-plan cadence $=S$). Offline pre-gates
confirmed fine-stride targets are well-posed (displacement $4.7\times$ the encoder noise
floor at $S{=}5$) and that the fine-scale CEM cost landscape \emph{sharpens} (true-plan
ranking percentile $0.93$ at $S{=}5$ vs.\ $0.82$ at $S{=}25$; cost dynamic range
$3\times$ larger).

\begin{table}[h]
\centering
\caption{Success rate vs.\ subgoal scale (controlled protocol). A clear interior optimum
at $S{=}10$: $+6.1/+10.1$pp over the paper configuration ($+8.7$pp at \degc{5} with 8
seeds, $\mathrm{SE}\approx0.9$pp, positive in every seed). $S{=}10$ pays almost no
penalty between the loose and strict criteria; $S{=}25$ drops ${\sim}5$pp.}
\label{tab:scale}
\begin{tabular}{lcccc}
\toprule
& $S{=}5$ & $S{=}10$ & $S{=}15$ & $S{=}25$ (paper config) \\
\midrule
SR \degc{20} (seeds 42--45)  & 92.3 & \ours{97.2} & 96.8 & 91.1 \\
SR \degc{5} (seeds 42--45)   & 91.9 & \ours{96.4} & 94.4 & 86.3 \\
SR \degc{5} (8 seeds, 42--49)& 92.9 & \ours{96.4} & 94.4 & 87.7 \\
Diffusion calls / episode    & 60   & 30          & 20   & 12   \\
\bottomrule
\end{tabular}
\end{table}

\paragraph{Mechanism: scale matching, not cadence.} A cadence-isolation control ---
re-anchor every 5 steps against \emph{unchanged} 25-step targets --- \emph{collapses} to
$78.6\%/77.8\%$ ($-12.5$pp below baseline): each re-draft re-samples a distant subgoal
from a wide conditional posterior, and CEM chases a churning target. Together with the
monotone commitment losses of \Cref{sec:neg}, the correct statement is: \textbf{closed-loop
success is governed by matching the control loop to the subgoal scale; finer matched
scales win, to an interior optimum} --- not ``re-plan more,'' not ``re-plan less.''

\paragraph{Prediction error is horizon-intrinsic.} Across all four drafters, offline
prediction error collapses onto the \emph{absolute} look-ahead horizon, independent of
block position ($+10$ steps: $0.105$ vs.\ $0.112$; $+15$: $0.132$ vs.\ $0.140$), and the
$S{=}5$ drafter's \emph{third} block position ($+15$, $0.140$) outpredicts the $S{=}25$
drafter's \emph{first} ($+25$, $0.159$). ``Suffix decay'' is task-future uncertainty, not
an architectural defect --- retiring intra-block architecture work.

% ============================================================
\section{Negative Results: The DSpark Port}
\label{sec:neg}

We implemented the published DSpark \citep{dspark} speculative-decoding composition on the GDM and
evaluated each component in the closed loop (controlled protocol, $S{=}25$).

\begin{table}[h]
\centering
\caption{DSpark components ported literally. Lower latent error does not mean better
control: the refiner \emph{halved} offline prediction error while losing 11pp SR ---
offline fidelity anti-correlates with closed-loop SR. The mechanism is conditional-spread
collapse: pulling drafts toward the conditional mean produces less-reachable targets.
This closes refinement, learned-confidence scheduling, and blind commitment as directions.}
\label{tab:dspark-neg}
\begin{tabular}{lccl}
\toprule
Arm & SR \degc{20} & SR \degc{5} & Verdict \\
\midrule
GDM, re-draft every step (reference)          & 91.9 & 83.1 & --- \\
+ refinement head (fixed-1, same cadence)     & 81.1 & 72.2 & refinement $-11$pp \\
+ blind commit-3                              & 73.8 & 50.0 & commitment $-18$pp \\
+ learned-confidence adaptive commit          & 64.5 & 52.0 & worse than blind \\
raw open-loop chain (commit-6)                & 72.6 & 63.7 & healthy chain still loses \\
\bottomrule
\end{tabular}
\end{table}

% ============================================================
\section{Random Initialization: The Paper's Hardest Setting}
\label{sec:randinit}

\begin{table}[h]
\centering
\caption{Random initialization, per-episode random goals (paper III-A protocol). The
protocol bug that once produced a spurious $100\%$ was fixed \emph{before} these runs;
the reproduced baseline scoring $78.1\%$ on the identical code path is the control.
$S{=}10$: \ours{$512/512$ episodes over two independent seeds, zero failures at both
tolerances}. The paper's criterion mapping is ambiguous for this cell (its $82.42$ aligns
with our strict-criterion baseline); under either reading, fine scale moves it to $100\%$
on every seed measured so far.}
\label{tab:randinit}
\begin{tabular}{lcc}
\toprule
& SR \degc{20} & SR \degc{5} \\
\midrule
LeWM (flat, goal image; paper)              & \multicolumn{2}{c}{0.00} \\
FF-JEPA Det (paper)                         & \multicolumn{2}{c}{81.25} \\
FF-JEPA GDM (paper)                         & \multicolumn{2}{c}{82.42} \\
FF-JEPA GDM (our faithful repro, $n{=}256$, seed 42) & 96.5 & 78.1 \\
Ours, $S{=}5$ ($n{=}256$, seed 42)          & 100.0 & 94.5 \\
Ours, $S{=}10$ ($n{=}256\times2$ seeds, 42--43) & \ours{100.0} & \ours{100.0} \\
\bottomrule
\end{tabular}
\end{table}

\paragraph{Data-completeness note.} A second $S{=}5$ seed and a third/fourth $S{=}10$
seed were queued (Isambard job 5578762) to extend this table to $n{=}256\times2$ ($S{=}5$)
and $n{=}256\times4$ ($S{=}10$) — matching the power of the rest of the paper's tables.
We could not locate a result for that job in any local log, results file, or the remote
cluster (unreachable at the time of writing) and do not report it. \emph{Every number
in this table is a directly verified success rate}; do not cite figures beyond what is
tabulated here (in particular, not ``1024 episodes'' or ``four seeds'') until the queued
arm is confirmed and this note is removed.

% ============================================================
\section{Reality-Verified Speculative Consumption}
\label{sec:spec}

The surviving DSpark idea is its \emph{skeleton}: draft a block cheaply, verify, accept
the surviving prefix. Our port replaces the learned verifier with reality: draft the
$N$-block once; after each hop, compare the \emph{achieved} latent to the subgoal just
pursued; within tolerance $\tau$ $\Rightarrow$ consume the next drafted position (no
diffusion call); otherwise $\Rightarrow$ re-draft from the achieved state. Re-anchoring is
never sacrificed when reality diverges. Offline calibration showed 8-step DDIM drafts are
indistinguishable from 50-step drafts (matched-noise distance $0.03$; spread preserved);
both transferred to the closed loop.

\begin{table}[h]
\centering
\caption{Speculative consumption (controlled protocol). NFE $=$ denoising function
evaluations per episode (calls $\times$ DDIM steps). At 8-seed power, verified coasting
is \degc{20}-neutral at $S{=}10$ but costs ${\sim}3.5$pp at \degc{5}; at $S{=}5$ the sign
\emph{reverses} (spec-accept beats its own every-step reference on both criteria) --- when
re-drafts churn a close target, skipping verified re-drafts is better control. The full
$S{=}10$ system stays $+6.3/+5.2$pp ahead of the paper configuration (8 seeds) at
$4.4\times$ lower diffusion cost.}
\label{tab:spec}
\begin{tabular}{lcccc}
\toprule
Arm ($S{=}10$ unless noted) & SR \degc{20} & SR \degc{5} & call ratio & NFE/ep \\
\midrule
every-step drafting, $k{=}50$ (reference)          & 97.2 & 96.4 & 1.00 & 1500 \\
spec-accept $\tau{=}0.15$, $k{=}8$                  & 98.0 & 95.2 & 0.67 & 161 \\
spec-accept $\tau{=}0.20$, $k{=}8$                  & \ours{98.0} & \ours{94.8} & \ours{0.56} & \ours{135} \\
spec-accept $\tau{=}0.30$, $k{=}8$                  & 96.8 & 94.8 & 0.45 & 108 \\
spec-accept $\tau{=}0.20$, $k{=}50$ (k-isolation)   & 98.0 & 95.6 & 0.55 & 825 \\
blind commit-2 (no verification)                   & 98.8 & 96.4 & 0.50 & 750 \\
\midrule
$S{=}5$ every-step (reference)                     & 92.3 & 91.9 & 1.00 & 3000 \\
$S{=}5$ spec-accept $\tau{=}0.20$, $k{=}8$          & \ours{98.0} & 94.8 & 0.42 & 202 \\
\midrule
\multicolumn{5}{l}{\emph{8-seed extension (seeds 42--49; $\mathrm{SE}\approx0.9$pp)}} \\
$S{=}10$ every-step                                 & 98.0 & \ours{96.4} & 1.00 & 1500 \\
$S{=}10$ spec-accept $\tau{=}0.20$, $k{=}8$         & 97.4 & 92.9 & 0.56 & \ours{135} \\
$S{=}5$ every-step                                  & 92.9 & 92.9 & 1.00 & 3000 \\
$S{=}5$ spec-accept $\tau{=}0.20$, $k{=}8$          & \ours{98.4} & \ours{96.2} & 0.42 & 202 \\
\midrule
paper config ($S{=}25$ every-step, 8 seeds \degc{5})& 91.1 & 87.7 & --- & 600 \\
\bottomrule
\end{tabular}
\end{table}

% ============================================================
\section{Is the Diffusion Drafter Necessary?}
\label{sec:regress}

A deterministic MLP regressor ($z_{cond}\to$ block, the offline-strong JointMLP class,
weight decay $10^{-3}$) as the subgoal source, re-anchoring every replan, matches the
diffusion drafter's \emph{offline} prediction error on the same windows ($0.30$ vs.\
$0.32$ at $S{=}10$). Closed loop, it loses everywhere.

\begin{table}[h]
\centering
\caption{Diffusion vs.\ deterministic regression as the subgoal source (seeds 42--45).
\textbf{The drafter's stochasticity is load-bearing}: at matched-or-better offline
fidelity, the zero-spread mean target loses everywhere, and its $-11.7$pp strict-criterion
loss at $S{=}25$ quantitatively mirrors the refinement head's $-11$pp (\Cref{tab:dspark-neg}),
closing the spread thesis causally --- any operator that regresses subgoals toward the
conditional mean damages control.}
\label{tab:regress}
\begin{tabular}{lcccc}
\toprule
& \multicolumn{2}{c}{$S{=}10$} & \multicolumn{2}{c}{$S{=}25$} \\
& \degc{20} & \degc{5} & \degc{20} & \degc{5} \\
\midrule
diffusion drafter (reference) & 97.2 & 96.4 & 91.1 & 86.3 \\
deterministic regressor       & 91.9 & 91.9 & 86.3 & 74.6 \\
\midrule
$\Delta$                      & $-5.2$ & $-4.4$ & $-4.8$ & $-11.7$ \\
\bottomrule
\end{tabular}
\end{table}

% ============================================================
\section{Failure Anatomy}
\label{sec:failure}

Instrumented traces (every re-plan's achieved latent and drafted subgoal) over all 8
seeds of $S{=}10$: \textbf{18 failures in 496 episodes, all tolerance-boundary
near-misses} --- final latent distance-to-goal within the success distribution's 95th
percentile (one failure sits \emph{closer} than the median success: a pure angle miss).
Zero unreachable-subgoal failures, zero divergence, zero stuck episodes; successful
episodes' commanded-to-achieved step ratio has median $1.06$. FF-JEPA's own reported
failure mode (``subgoal looks right, agent cannot reach it'') is \emph{eliminated} at fine
scale; the residue is quantization at the success threshold.

% ============================================================
\section{Horizon Scaling (Fixed Population)}
\label{sec:horizon}

Goal offsets $t\in\{25,50,75,100\}$ on the \emph{identical} 256 episodes at every offset
(reindexed from the $t{=}150$ set), budget $2t$, seeds 42/43 --- so horizon is the only
variable, free of the eligibility confound that invalidates independent per-offset draws.

\begin{table}[h]
\centering
\small
\setlength{\tabcolsep}{4.5pt}
\caption{Fixed-population horizon sweep ($n{=}512$/cell). \textbf{The baseline decays
monotonically at the strict criterion ($-11.9$pp from $t{=}25$ to $t{=}100$) while
$S{=}10$ is flat from $t{=}50$ onward --- the gap widens monotonically.} Spec-accept tracks
every-step within $1.2$pp across the curve (call ratio ${\approx}0.57$ at long horizons,
${\approx}0.85$ at $t{=}25$: the cost win is a long-horizon phenomenon).}
\label{tab:horizon}
\begin{tabular}{lcccccc}
\toprule
& \multicolumn{2}{c}{$S{=}25$ (FF-JEPA)} & \multicolumn{2}{c}{$S{=}10$ every-step}
& \multicolumn{2}{c}{$S{=}10$ $+$ spec-accept} \\
\cmidrule(lr){2-3}\cmidrule(lr){4-5}\cmidrule(lr){6-7}
$t$ & \degc{20} & \degc{5} & \degc{20} & \degc{5} & \degc{20} & \degc{5} \\
\midrule
25  & 98.1 & 93.8 & \ours{99.8} & \ours{97.1} & 99.6 & 95.1 \\
50  & 93.6 & 85.0 & \ours{97.5} & \ours{93.8} & 97.1 & 93.9 \\
75  & 92.6 & 83.2 & \ours{98.2} & \ours{93.8} & 97.9 & 93.0 \\
100 & 93.8 & 81.8 & 97.9 & \ours{93.8} & \ours{98.6} & 92.6 \\
\midrule
\degc{5} gap ($S{=}10-$ baseline) & & & \multicolumn{2}{c}{$+3.3\to+8.8\to+10.5\to+11.9$} & & \\
\bottomrule
\end{tabular}
\end{table}

\paragraph{Cross-paper regime (VLWM fixed-budget-50).} We additionally evaluate under the
VLWM \citep{vlwm} fixed-budget-50 regime --- a harsher protocol than FF-JEPA's
final-window budget of $2t$, holding the step budget at 50 for \emph{every} offset
(\Cref{tab:vlwm}).

\begin{table}[h]
\centering
\caption{Cross-paper cell: our three arms under VLWM's fixed-budget-50 regime, all
\degc{20}, seed means. Rightmost columns are the published PushT numbers of VLWM and of
LeWM (flat) under that regime. Budget starvation reproduces the published collapse, yet
within the identical budget our arms sit ${\sim}3\times$ above VLWM's bars. Curve
\emph{shapes} are comparable across papers; absolute numbers are not (start/goal sampling
differs), so we cite this as ``same budget regime,'' not a like-for-like win.
$^{\dagger}$At $t{=}25$ the fixed-50 budget equals FF-JEPA's native $2t$ budget, so the two
protocols coincide (the row is the $t{=}25$ cell of \Cref{tab:horizon}); the budget binds
only for $t\geq50$, which is where the published VLWM/LeWM bars begin.}
\label{tab:vlwm}
\begin{tabular}{lccc cc}
\toprule
& \multicolumn{3}{c}{Ours (fixed-budget-50)} & \multicolumn{2}{c}{Published} \\
\cmidrule(lr){2-4}\cmidrule(lr){5-6}
$t$ & $S{=}25$ & $S{=}10$ & $S{=}10$ $+$ spec & VLWM & LeWM (flat) \\
\midrule
25$^{\dagger}$ & 98.1 & \ours{99.8} & 99.6 & 94 & 90 \\
50  & 83.6 & 88.7 & \ours{90.0} & 60 & 46 \\
75  & 55.7 & 59.8 & \ours{60.9} & 20 & 12 \\
100 & 23.4 & 24.4 & \ours{27.3} & 12 & 8  \\
\bottomrule
\end{tabular}
\end{table}

Four read-offs, all consistent with the main story: (a) \textbf{budget starvation
reproduces the collapse} --- confirming VLWM's PushT cliff is largely a budget artifact,
which validates our protocol note; (b) \textbf{within their regime we sit ${\sim}3\times$
above their bars} ($60.9$ vs.\ $20$ at $t{=}75$), citable as ``same budget regime'' with
the start/goal-sampling caveat; (c) \textbf{spec-accept beats every-step at every offset}
under time pressure --- the anti-churn effect (\Cref{sec:spec}) generalizes to $S{=}10$
once budget binds, a new data point for the $\tau$-trade; (d) \textbf{the $S{=}10{-}S{=}25$
gap compresses under starvation} --- scale matching buys precision, not raw speed.

% ============================================================
\section{Headline Comparison}
\label{sec:headline}

\begin{table}[h]
\centering
\small
\setlength{\tabcolsep}{5pt}
\caption{Where we stand, both criteria per cell (\degc{20} $=$ the criterion enforced by
the benchmark's shipped code; \degc{5} $=$ the strict precision criterion its prose
states). ``Ours'' modifies FF-JEPA's GDM system only in subgoal scale, loop matching, and
consumption policy --- encoder, drafter architecture, training recipe, and CEM are
unchanged. $^{p}$The paper publishes one number per cell without stating its criterion;
each is placed on the side it empirically aligns with, with $\cdot$ marking the untested
side. Comparisons are valid only within a side.}
\label{tab:headline}
\begin{tabular}{lccccc}
\toprule
& & \multicolumn{2}{c}{FF-JEPA GDM} & \multicolumn{2}{c}{\ours{Ours}} \\
\cmidrule(lr){3-4}\cmidrule(lr){5-6}
Setting (SR \degc{20} / \degc{5}) & LeWM (flat) & published$^{p}$ & our repro & $S{=}10$ & $+$ spec-accept \\
\midrule
Short horizon, $t{=}25$          & 94.5 & 96.1 / $\cdot$ & 95.3 / $\cdot$ & \ours{98.4} / 94.1 & 98.4 / \ours{95.3} \\
Long horizon, $t{=}75$           & 3.5  & 91.8 / $\cdot$ & 90.0 / 78.8 & 97.1 / \ours{93.8} & \ours{97.9} / 91.6 \\
Long horizon, $t{=}150$$^{a}$    & ---  & --- & 91.1 / 87.7 & \ours{98.0 / 96.4} & 97.4 / 92.9 \\
Horizon sweep, $t{=}100$$^{b}$   & ---  & --- & 93.8 / 81.8 & 97.9 / \ours{93.8} & \ours{98.6} / 92.6 \\
Random initialization$^{c}$      & 0.0  & $\cdot$ / 82.4 & 96.5 / 78.1 & \ours{100.0 / 100.0}$^{d}$ & --- \\
\midrule
Diffusion NFE / episode$^{a}$    & 0    & --- & 600 & 1500 & \ours{135} \\
\bottomrule
\end{tabular}

\vspace{1mm}
{\footnotesize $^{a}$Controlled protocol, 8 seeds. $^{b}$Fixed-population sweep,
$n{=}512$/cell. $^{c}$$n{=}512$ over 2 seeds, zero failures at both tolerances;
spec-accept not run in this mode (the drafter is the tested component).
$^{d}$Two further seeds are queued but unconfirmed --- see the data-completeness
note in \Cref{sec:randinit}; do not extrapolate beyond the verified $n{=}512$.}
\end{table}

% ============================================================
\section{Generalization to Further Environments}
\label{sec:generalize}
% All results above are on PushT. The following environments extend the
% subgoal-scale + spec-accept claims beyond a single task. Tables are
% placeholders until the runs land.

Every result above is on PushT. To establish that the subgoal-scale law and
reality-verified speculative consumption are properties of hierarchical latent planning
rather than of one task, we extend to three environments of increasing structural
distance from PushT. \emph{These runs are in progress; tables below are placeholders.}

\subsection{Reacher}
\label{sec:reacher}
A low-dimensional continuous-control reach task --- the recon that decides whether the
scale law transfers to a non-pushing morphology. Same protocol: retrain the drafter at
$S\in\{5,10,15,25\}$, report SR vs.\ scale and the spec-accept efficiency arm.
\tbdblock{Reacher: SR vs.\ subgoal scale + spec-accept NFE. Recon in progress.}

\subsection{TwoRoom}
\label{sec:tworoom}
A navigation task with a bottleneck between rooms, probing whether scale matching holds
when the optimal subgoal density is spatially non-uniform.
\tbdblock{TwoRoom: SR vs.\ subgoal scale under goal-conditioned planning.}

\subsection{OGBench-Cube}
\label{sec:ogbench}
A manipulation benchmark \citep{ogbench} for the strongest generalization claim: multi-object
state and longer effective horizons, where the ${\sim}11\times$ diffusion-cost reduction
matters most.
\tbdblock{OGBench-Cube: headline SR / NFE vs.\ FF-JEPA baseline.}

% ============================================================
\section{Discussion and Limitations}
\label{sec:discussion}

\paragraph{What the story is.} Closed-loop success in this planner is governed by matching
the control loop to the subgoal scale, and the drafter's conditional \emph{spread} is
load-bearing: every operator we tried that collapses drafts toward the conditional mean
(refinement, blind commitment, deterministic regression) damaged control despite equal or
better offline fidelity. The efficiency win is orthogonal --- reality is a free verifier.

\paragraph{Limitations.} (i) FF-JEPA is a preliminary paper with no released code; our
``faithful'' configuration is a principled reconstruction where the paper is silent, not a
verified match on every axis. (ii) All confirmed results are on PushT; cross-environment
generality (\Cref{sec:generalize}) is \TBD --- Reacher and TwoRoom drivers are implemented
and committed but not yet run (pending cluster availability); OGBench-Cube is not yet
started. (iii) The \degc{20}/\degc{5} criterion mapping for the paper's own cells is
inferred from empirical alignment, not stated by the paper; we mark every cross-paper
comparison with the side it is valid on. (iv) The $t{=}150$ cell is a beyond-paper probe,
not a reproduction claim. (v) The random-initialization table (\Cref{sec:randinit}) is
one arm short of its intended power --- a queued confirmation run ($S{=}5$ second seed,
$S{=}10$ third/fourth seeds) has no recoverable result and is reported as such rather
than assumed.

% ============================================================
\section{Conclusion}
\label{sec:conclusion}
\tbdblock{This is a placeholder conclusion for a preliminary, PushT-only draft --- it
will be rewritten once \Cref{sec:generalize} has real numbers and is not meant to stand
as this paper's final framing.}

On PushT, subgoal scale matched to the control loop is the strongest lever we found in
FF-JEPA-style hierarchical latent planning: a fine matched stride ($S{=}10$) beats the
paper configuration across every setting we measured and, on the paper's own protocols,
exceeds its published numbers. Reality-verified speculative consumption recovers
${\sim}11\times$ of the added diffusion cost at neutral success rate. Whether either
claim is a property of hierarchical latent planning in general, versus an artifact of
PushT specifically, is exactly what \Cref{sec:generalize} is for and is not yet
answered --- that is the remaining work, not a formality.

% ============================================================
\bibliographystyle{plainnat}
\bibliography{references}
% references.bib currently holds stub/placeholder entries for
% ffjepa, lewm, cem, vlwm, ogbench, dspark -- fill with real
% metadata before this leaves preliminary-draft status.

\end{document}
